% Anti-linear completion of Hessian 216 and the affine-hull null cone.
\subsection*{Anti-linear conjugation completes the Hessian group and exposes a four-ray finite null cone}

The projective one-qutrit Clifford/Hessian symmetry is the already-certified
central quotient
\[
 648/C_3=216,\qquad G_{216}\cong
 \mathbb F_3^2:\SL(2,3).
\]
In the frozen normal form $Z^aX^b\omega^c$, coefficient conjugation acts as
\[
 (a,b,c)\longmapsto(-a,b,-c).
\]
Consequently its action on $H_{27}/Z(H_{27})\cong\mathbb F_3^2$ is the
determinant-minus-one reflection $\operatorname{diag}(-1,1)$.  Adjoining this
reflection gives
\[
 \mathbb F_3^2:\GL(2,3)=\AGL(2,3),\qquad |\AGL(2,3)|=432,
\]
exactly the affine factor independently recovered from the ramified
$[45,12,6]_3$ VM code.

On the four affine direction classes
\[
 (1,0),\ (0,1),\ (1,1),\ (1,2),
\]
the determinant-one subgroup has image $A_4$, while conjugation supplies the
odd transposition exchanging the last two classes.  The extended image is
therefore all of $S_4$.

Independently, the nine-dimensional hull has a three-dimensional nondegenerate
quotient over $\mathbb F_3$ with Gram matrix
\[
 Q=\begin{pmatrix}0&1&1\\1&0&1\\1&1&0\end{pmatrix}.
\]
Its four projective isotropic rays are
\[
 [1,0,0],\ [0,1,0],\ [0,0,1],\ [1,1,1].
\]
Direct exhaustion gives $|O(Q)|=48$ and $|SO(Q)|=24$; the special orthogonal
group acts faithfully as $S_4$ on those four rays.  Under the dictionary
\[
 (1,0)\leftrightarrow[1,0,0],\quad
 (0,1)\leftrightarrow[0,1,0],\quad
 (1,1)\leftrightarrow[0,0,1],\quad
 (1,2)\leftrightarrow[1,1,1],
\]
its ray permutations equal the full $\GL(2,3)$ direction permutations.

Thus the four Hesse parallel classes are exactly the projective isotropic set
of the affine-hull quotient, and the missing factor two from $216$ to $432$ is
realized by anti-linear qutrit orientation reversal.

This is a finite orthogonal geometry over $\mathbb F_3$.  It does not supply a
continuum Lorentz metric, causal cone, CPT theorem, or measured CP observable.
